\documentclass[11pt,paper=a4,answers]{exam}
\usepackage{graphicx,lastpage}
\usepackage{upgreek}
\usepackage{censor}
\usepackage{gensymb}
\usepackage{amsmath}
\usepackage{multicol}
\usepackage{enumitem}
\usepackage{tasks}
\usepackage{adjustbox}
\usepackage{xcolor}
\usepackage{tfrupee}
\setlength{\voffset}{0.25in}
\usepackage{tabularx}
\newcolumntype{Y}{>{\centering\arraybackslash}X}
%==============================================================
% Customise Non-Essential Packages Below
% =============================================================
\usepackage[most]{tcolorbox}
\usepackage{eso-pic}
\usepackage{tikz}
\AddToShipoutPictureBG{%
\begin{tikzpicture}[remember picture,overlay]
\foreach \x in {-8,-4,0,4,8,12,16,20}
\foreach \y in {-8,-4,0,4,8,12,16,20} {
\node[opacity=0.05,rotate=45,anchor=center] at ([xshift=\x cm,yshift=\y cm]current page.center) {%
\begin{minipage}{2cm}
\centering
\includegraphics[width=2cm]{images/logo_black.png}\\
\small preppeo.com
\end{minipage}%
};
}
\end{tikzpicture}%
}
\printanswers

%==============================================================
% Customise Non-Essential Packages Above
% =============================================================
\definecolor{svgray}{rgb}{0.90, 0.90, 0.90}
\graphicspath{{images/}}

\usepackage[normalem]{ulem}
\renewcommand\ULthickness{1pt}   %%---> For changing thickness of underline
\setlength\ULdepth{3.5ex}%\maxdimen ---> For changing depth of underline
\renewcommand{\baselinestretch}{1}
\chead{
\includegraphics[scale=0.1,height=2cm ]{logo.png}
}
\pagestyle{empty}

\pagestyle{headandfoot}
\headrule
\cfoot{W: https://preppeo.com; M: +91-8130711689}

\begin{document}
%==============================================================
% Customise Question paper details below this
% =============================================================
\begin{center}
    \centering
    \renewcommand{\arraystretch}{1.5}
    \begin{tabularx}{\textwidth}{YYY}
    & \normalsize\bf{{\Large{{Quantitative Reasoning}}}} & \\
    By Preppeo & \normalsize\bf{{gre\_lid\_016}} & SAT\\
    \normalsize{{\textbf{{Unit:}} Advanced Math}} & \normalsize{{\textbf{{Chapter:}} Systems of equations in two variables}} & \normalsize{{\textbf{{Topic:}} Linear-Quadratic Systems}} \\
    \end{tabularx}
\end{center}
\par\noindent\rule{\textwidth}{1pt}
% =============================================================
% Customise Question paper details above this
% =============================================================

\begin{questions}
\pointsinrightmargin
\pointsdroppedatright
\marksnotpoints
\pointpoints{mark}{marks}
\pointformat{\boldmath\themarginpoints}
\bracketedpoints

% =============================================================
% Question Begin
% =============================================================
% Question 1
% Category: Practice | Difficulty: Medium | Question Type: Multiple Choice
\question
A system of equations consists of the parabola $y = x^2 - 8x + 19$ and the horizontal line $y = 3$. At how many points do the graphs of these two equations intersect in the $xy$-plane?
\begin{tasks}(4)
    \task Zero
    \task One
    \task Two
    \task Infinitely many
\end{tasks}

\begin{solution}
\textbf{Conceptual Explanation:} \\
The intersection points represent the real solutions to the system. We can determine the number of solutions by setting the equations equal to each other and analyzing the discriminant ($b^2 - 4ac$) of the resulting quadratic.

\vspace{20pt}

\textbf{Calculation and Logic:} \\
Set the equations equal to each other: \\
$x^2 - 8x + 19 = 3$

\vspace{10pt}

Move all terms to one side to set the equation to zero: \\
$x^2 - 8x + 16 = 0$

\vspace{10pt}

Identify the coefficients: $a = 1, b = -8, c = 16$. \\
Calculate the discriminant: \\
$D = (-8)^2 - 4(1)(16)$ \\
$D = 64 - 64 = 0$

\vspace{10pt}

Since the discriminant is exactly zero, the system has only one real solution, meaning the line is tangent to the parabola at its vertex.

\vspace{25pt}

Answer: \colorbox{yellow!100}{B}
\end{solution}

\vspace{40pt}

% Question 2
% Category: Practice | Difficulty: Hard | Question Type: Multiple Choice
\question
If $(x, y)$ is a solution to the system of equations below and $x > 0$, what is the value of $x$?
\[ \begin{cases} y = x^2 - 12 \\ y = x \end{cases} \]
\begin{tasks}(4)
    \task 3
    \task 4
    \task 6
    \task 12
\end{tasks}

\begin{solution}
\textbf{Conceptual Explanation:} \\
To find the coordinates of the intersection points, substitute the expression for $y$ from the linear equation into the quadratic equation. Then, solve the quadratic and apply the given constraint for $x$.

\vspace{20pt}

\textbf{Calculation and Logic:} \\
Substitute $x$ for $y$ in the first equation: \\
$x = x^2 - 12$

\vspace{10pt}

Rearrange into standard form: \\
$x^2 - x - 12 = 0$

\vspace{10pt}

Factor the quadratic: \\
$(x - 4)(x + 3) = 0$

\vspace{10pt}

Solve for $x$: \\
$x = 4$ or $x = -3$

\vspace{10pt}

Because the problem specifies that $x > 0$, the only valid value is 4.

\vspace{25pt}

Answer: \colorbox{yellow!100}{B}
\end{solution}

\vspace{40pt}

% Question 3
% Category: Practice | Difficulty: Medium | Question Type: Multiple Choice
\question
The graph of a system of equations consists of a parabola and a line. The parabola is defined by $y = x^2 - 5$ and the line is defined by $y = 4$. What are the $x$-coordinates of the points where the line and the parabola intersect?
\begin{tasks}(2)
    \task $x = 3$ and $x = -3$
    \task $x = 2$ and $x = -2$
    \task $x = 1$ and $x = -1$
    \task $x = 5$ and $x = -5$
\end{tasks}

\begin{solution}
\textbf{Conceptual Explanation:} \\
The intersection occurs when both equations have the same $y$-value. We set the quadratic part of the parabola equal to the constant value of the horizontal line.

\vspace{20pt}

\textbf{Calculation and Logic:} \\
Set the equations equal: \\
$x^2 - 5 = 4$

\vspace{10pt}

Isolate the squared variable: \\
$x^2 = 9$

\vspace{10pt}

Take the square root of both sides: \\
$x = \pm 3$

\vspace{10pt}

The intersection points occur at $x = 3$ and $x = -3$.

\vspace{25pt}

Answer: \colorbox{yellow!100}{A}
\end{solution}

\vspace{40pt}

% Question 4
% Category: Practice | Difficulty: Hard | Question Type: Multiple Choice
\question
A system of equations is given by $y = x^2 + 10x + c$ and $y = -5$. If the system has exactly one real solution, what is the value of the constant $c$?
\begin{tasks}(4)
    \task 20
    \task 25
    \task 30
    \task -20
\end{tasks}

\begin{solution}
\textbf{Conceptual Explanation:} \\
For a linear-quadratic system to have exactly one solution, the line must be tangent to the parabola. Algebraically, this means the combined quadratic equation must have a discriminant equal to zero.

\vspace{20pt}

\textbf{Calculation and Logic:} \\
Set the equations equal: \\
$x^2 + 10x + c = -5$

\vspace{10pt}

Rearrange into standard form ($ax^2 + bx + c = 0$): \\
$x^2 + 10x + (c + 5) = 0$

\vspace{10pt}

Identify the coefficients: $a = 1, b = 10, \text{constant} = c + 5$. \\
Set the discriminant to zero: \\
$(10)^2 - 4(1)(c + 5) = 0$

\vspace{10pt}

Solve for $c$: \\
$100 - 4c - 20 = 0$ \\
$80 = 4c$ \\
$c = 20$

\vspace{25pt}

Answer: \colorbox{yellow!100}{A}
\end{solution}

\vspace{40pt}

% Question 5
% Category: Practice | Difficulty: Hard | Question Type: Multiple Choice
\question
In the $xy$-plane, the line $y = 2x + 1$ intersects the parabola $y = x^2 - x + 3$ at two points. What is the sum of the $x$-coordinates of these two points?
\begin{tasks}(4)
    \task 1
    \task 2
    \task 3
    \task 4
\end{tasks}

\begin{solution}
\textbf{Conceptual Explanation:} \\
The sum of the $x$-coordinates of the intersection points is the sum of the roots of the quadratic equation formed by setting the two functions equal. We can use the formula $-b/a$ for the sum of the roots.

\vspace{20pt}

\textbf{Calculation and Logic:} \\
Set the equations equal: \\
$x^2 - x + 3 = 2x + 1$

\vspace{10pt}

Move all terms to one side: \\
$x^2 - 3x + 2 = 0$

\vspace{10pt}

Identify the coefficients: $a = 1, b = -3$. \\
Use the sum of roots formula: \\
Sum $= -\frac{b}{a} = -\frac{-3}{1} = 3$

\vspace{25pt}

Answer: \colorbox{yellow!100}{C}
\end{solution}

% --- PART 1 continued: Linear-Quadratic Systems (6-25) ---

% Question 6
% Category: Practice | Difficulty: Medium | Question Type: Multiple Choice
\question
Which of the following points $(x, y)$ is a solution to the system of equations below?
\[ \begin{cases} y = x^2 - 10 \\ y = 3x \end{cases} \]
\begin{tasks}(4)
    \task $(2, 6)$
    \task $(5, 15)$
    \task $(-2, 6)$
    \task $(5, 25)$
\end{tasks}

\begin{solution}
\textbf{Conceptual Explanation:} \\
To find the solution to the system, substitute the linear expression for $y$ into the quadratic equation. Solve for $x$, and then use either equation to find the corresponding $y$-value.

\vspace{20pt}

\textbf{Calculation and Logic:} \\
Set the equations equal: \\
$x^2 - 10 = 3x$

\vspace{10pt}

Rearrange to standard form: \\
$x^2 - 3x - 10 = 0$

\vspace{10pt}

Factor the quadratic: \\
$(x - 5)(x + 2) = 0$ \\
$x = 5$ or $x = -2$

\vspace{10pt}

Find the $y$-value for $x = 5$ using $y = 3x$: \\
$y = 3(5) = 15$ \\
The point is $(5, 15)$.

\vspace{25pt}

Answer: \colorbox{yellow!100}{B}
\end{solution}

\vspace{40pt}

% Question 2
% Category: Practice | Difficulty: Hard | Question Type: Multiple Choice
\question
The line $y = kx - 4$ is tangent to the parabola $y = x^2$ at exactly one point. If $k > 0$, what is the value of $k$?
\begin{tasks}(4)
    \task 2
    \task 4
    \task 8
    \task 16
\end{tasks}

\begin{solution}
\textbf{Conceptual Explanation:} \\
A line is tangent to a parabola when the system has exactly one real solution. This means the discriminant of the resulting quadratic must be equal to zero.

\vspace{20pt}

\textbf{Calculation and Logic:} \\
Set the equations equal: \\
$x^2 = kx - 4$

\vspace{10pt}

Rearrange to standard form: \\
$x^2 - kx + 4 = 0$

\vspace{10pt}

Identify coefficients: $a = 1, b = -k, c = 4$. \\
Set the discriminant ($b^2 - 4ac$) to zero: \\
$(-k)^2 - 4(1)(4) = 0$ \\
$k^2 - 16 = 0$

\vspace{10pt}

Solve for $k$: \\
$k = 4$ or $k = -4$

\vspace{10pt}

Since the problem states $k > 0$, the value is 4.

\vspace{25pt}

Answer: \colorbox{yellow!100}{B}
\end{solution}

\vspace{40pt}

% Question 8
% Category: Practice | Difficulty: Medium | Question Type: Multiple Choice
\question
How many real solutions does the following system of equations have?
\[ \begin{cases} y = x^2 + 4x + 6 \\ y = -2 \end{cases} \]
\begin{tasks}(4)
    \task Zero
    \task One
    \task Two
    \task Infinitely many
\end{tasks}

\begin{solution}
\textbf{Conceptual Explanation:} \\
The number of solutions can be found by seeing if the quadratic expression can ever equal the constant value. We can also consider the vertex of the parabola.

\vspace{20pt}

\textbf{Calculation and Logic:} \\
Find the vertex of $y = x^2 + 4x + 6$: \\
$x = -b/2a = -4/2 = -2$ \\
$y = (-2)^2 + 4(-2) + 6 = 4 - 8 + 6 = 2$

\vspace{10pt}

The vertex is $(-2, 2)$ and the parabola opens upward. \\
Since the minimum $y$-value is 2, the graph will never intersect the line $y = -2$.

\vspace{25pt}

Answer: \colorbox{yellow!100}{A}
\end{solution}

\vspace{40pt}

% Question 9
% Category: Practice | Difficulty: Hard | Question Type: Multiple Choice
\question
In the $xy$-plane, the parabola $y = x^2 - k$ and the line $y = 2$ intersect at two points. If the distance between these two points is 6, what is the value of $k$?
\begin{tasks}(4)
    \task 4
    \task 7
    \task 9
    \task 11
\end{tasks}

\begin{solution}
\textbf{Conceptual Explanation:} \\
The intersection points have the same $y$-coordinate (2). The distance between them is the difference between their $x$-coordinates.

\vspace{20pt}

\textbf{Calculation and Logic:} \\
Set the equations equal: \\
$x^2 - k = 2$ \\
$x^2 = k + 2$ \\
$x = \pm \sqrt{k + 2}$

\vspace{10pt}

The distance between $x_1$ and $x_2$ is: \\
$\sqrt{k + 2} - (-\sqrt{k + 2}) = 6$ \\
$2\sqrt{k + 2} = 6$ \\
$\sqrt{k + 2} = 3$

\vspace{10pt}

Square both sides: \\
$k + 2 = 9$ \\
$k = 7$

\vspace{25pt}

Answer: \colorbox{yellow!100}{B}
\end{solution}

\vspace{40pt}

% Question 10
% Category: Practice | Difficulty: Medium | Question Type: Multiple Choice
\question
Which of the following values of $x$ is a solution to the system below?
\[ \begin{cases} y = x^2 - x \\ y = 2x + 4 \end{cases} \]
\begin{tasks}(4)
    \task 1
    \task 2
    \task 4
    \task 6
\end{tasks}

\begin{solution}
\textbf{Conceptual Explanation:} \\
Set the quadratic and linear expressions equal to find the $x$-values where the graphs cross.

\vspace{20pt}

\textbf{Calculation and Logic:} \\
$x^2 - x = 2x + 4$ \\
$x^2 - 3x - 4 = 0$ \\
$(x - 4)(x + 1) = 0$

\vspace{10pt}

The solutions for $x$ are 4 and $-1$. \\
Looking at the choices, 4 is the correct answer.

\vspace{25pt}

Answer: \colorbox{yellow!100}{C}
\end{solution}

% --- PART 1 continued: Linear-Quadratic Systems (11-25) ---

% Question 11
% Category: Practice | Difficulty: Medium | Question Type: Multiple Choice
\question
A system of equations consists of $y = 2x^2 - 8x + 9$ and $y = 1$. How many solutions $(x, y)$ does this system have?
\begin{tasks}(4)
    \task Zero
    \task One
    \task Two
    \task Infinitely many
\end{tasks}

\begin{solution}
\textbf{Conceptual Explanation:} \\
The number of solutions is determined by the discriminant of the quadratic equation formed by setting the two expressions equal to each other.

\vspace{20pt}

\textbf{Calculation and Logic:} \\
Set the equations equal: \\
$2x^2 - 8x + 9 = 1$

\vspace{10pt}

Subtract 1 from both sides: \\
$2x^2 - 8x + 8 = 0$

\vspace{10pt}

Divide by 2 to simplify: \\
$x^2 - 4x + 4 = 0$

\vspace{10pt}

Calculate the discriminant ($b^2 - 4ac$): \\
$(-4)^2 - 4(1)(4) = 16 - 16 = 0$

\vspace{10pt}

Since the discriminant is 0, the system has exactly one solution.

\vspace{25pt}

Answer: \colorbox{yellow!100}{B}
\end{solution}

\vspace{40pt}

% Question 12
% Category: Practice | Difficulty: Hard | Question Type: Multiple Choice
\question
The line $y = 3x + c$ is tangent to the parabola $y = x^2 - x + 5$ at exactly one point. What is the value of $c$?
\begin{tasks}(4)
    \task 1
    \task 2
    \task 3
    \task 4
\end{tasks}

\begin{solution}
\textbf{Conceptual Explanation:} \\
For a line to be tangent to a parabola, the system must have exactly one solution. This requires the discriminant of the combined quadratic equation to be zero.

\vspace{20pt}

\textbf{Calculation and Logic:} \\
Set the equations equal: \\
$x^2 - x + 5 = 3x + c$

\vspace{10pt}

Move all terms to one side: \\
$x^2 - 4x + (5 - c) = 0$

\vspace{10pt}

Identify coefficients: $a = 1, b = -4, c = (5 - c)$. \\
Set the discriminant to zero: \\
$(-4)^2 - 4(1)(5 - c) = 0$ \\
$16 - 20 + 4c = 0$

\vspace{10pt}

Solve for $c$: \\
$-4 + 4c = 0$ \\
$4c = 4 \Rightarrow c = 1$

\vspace{25pt}

Answer: \colorbox{yellow!100}{A}
\end{solution}

\vspace{40pt}

% Question 13
% Category: Practice | Difficulty: Medium | Question Type: Multiple Choice
\question
At which of the following points does the line $y = x + 2$ intersect the parabola $y = x^2$?
\begin{tasks}(4)
    \task $(1, 3)$
    \task $(2, 4)$
    \task $(0, 2)$
    \task $(-1, 2)$
\end{tasks}

\begin{solution}
\textbf{Conceptual Explanation:} \\
Intersection points are found by setting the equations equal to each other and solving for $x$. Then, substitute those $x$-values into either equation to find $y$.

\vspace{20pt}

\textbf{Calculation and Logic:} \\
Set the equations equal: \\
$x^2 = x + 2$

\vspace{10pt}

Rearrange into standard form: \\
$x^2 - x - 2 = 0$

\vspace{10pt}

Factor the quadratic: \\
$(x - 2)(x + 1) = 0$ \\
$x = 2$ or $x = -1$

\vspace{10pt}

Find the $y$-value for $x = 2$: \\
$y = (2)^2 = 4$ \\
The intersection point is $(2, 4)$.

\vspace{25pt}

Answer: \colorbox{yellow!100}{B}
\end{solution}

\vspace{40pt}

% Question 14
% Category: Practice | Difficulty: Hard | Question Type: Multiple Choice
\question
If $(x, y)$ is a solution to the system below, what is a possible value of $x + y$?
\[ \begin{cases} y = x^2 - 7 \\ y = 6x \end{cases} \]
\begin{tasks}(4)
    \task -6
    \task 0
    \task 42
    \task 49
\end{tasks}

\begin{solution}
\textbf{Conceptual Explanation:} \\
First solve for $x$ by substitution, then find the corresponding $y$, and finally add them together.

\vspace{20pt}

\textbf{Calculation and Logic:} \\
$x^2 - 7 = 6x \Rightarrow x^2 - 6x - 7 = 0$ \\
$(x - 7)(x + 1) = 0$ \\
$x = 7$ or $x = -1$

\vspace{10pt}

Case 1 ($x = 7$): \\
$y = 6(7) = 42$ \\
Sum: $7 + 42 = 49$

\vspace{10pt}

Case 2 ($x = -1$): \\
$y = 6(-1) = -6$ \\
Sum: $-1 + (-6) = -7$

\vspace{10pt}

Choice D (49) is the only option that matches our results.

\vspace{25pt}

Answer: \colorbox{yellow!100}{D}
\end{solution}

\vspace{40pt}

% Question 15
% Category: Practice | Difficulty: Medium | Question Type: Multiple Choice
\question
How many points of intersection exist between $y = -x^2 + 5$ and $y = 10$?
\begin{tasks}(4)
    \task Zero
    \task One
    \task Two
    \task Infinitely many
\end{tasks}

\begin{solution}
\textbf{Conceptual Explanation:} \\
Compare the maximum value of the downward-opening parabola to the height of the horizontal line.

\vspace{20pt}

\textbf{Calculation and Logic:} \\
The parabola $y = -x^2 + 5$ has its vertex at $(0, 5)$ and opens downward. \\
The maximum $y$-value of this parabola is 5. \\
The line $y = 10$ is strictly above the maximum point of the parabola. \\
Therefore, there are no points of intersection.

\vspace{25pt}

Answer: \colorbox{yellow!100}{A}
\end{solution}

% Question 16
% Category: Practice | Difficulty: Hard | Question Type: Multiple Choice
\question
The line $y = mx$ is tangent to the parabola $y = x^2 + 9$ at exactly one point. If $m > 0$, what is the value of $m$?
\begin{tasks}(4)
    \task 3
    \task 6
    \task 9
    \task 18
\end{tasks}

\begin{solution}
\textbf{Conceptual Explanation:} \\
When a line passing through the origin is tangent to a parabola, the resulting quadratic equation must have a discriminant of zero.

\vspace{20pt}

\textbf{Calculation and Logic:} \\
Set the equations equal: \\
$x^2 + 9 = mx$ \\
$x^2 - mx + 9 = 0$

\vspace{10pt}

Identify coefficients: $a = 1, b = -m, c = 9$. \\
Set the discriminant to zero: \\
$(-m)^2 - 4(1)(9) = 0$ \\
$m^2 - 36 = 0$ \\
$m^2 = 36 \Rightarrow m = 6$ (since $m > 0$)

\vspace{25pt}

Answer: \colorbox{yellow!100}{B}
\end{solution}

\vspace{40pt}

% Question 17
% Category: Practice | Difficulty: Medium | Question Type: Multiple Choice
\question
If the line $x = 3$ intersects the parabola $y = x^2 - 2x + 5$ at point $P$, what is the $y$-coordinate of $P$?
\begin{tasks}(4)
    \task 5
    \task 8
    \task 11
    \task 14
\end{tasks}

\begin{solution}
\textbf{Conceptual Explanation:} \\
A vertical line $x = k$ intersects a function at the point $(k, f(k))$. Simply substitute the $x$-value into the quadratic equation.

\vspace{20pt}

\textbf{Calculation and Logic:} \\
Substitute $x = 3$ into the equation: \\
$y = (3)^2 - 2(3) + 5$ \\
$y = 9 - 6 + 5$ \\
$y = 8$

\vspace{25pt}

Answer: \colorbox{yellow!100}{B}
\end{solution}

\vspace{40pt}

% Question 18
% Category: Practice | Difficulty: Hard | Question Type: Multiple Choice
\question
The system below has exactly one solution $(x, y)$. What is the value of $x$?
\[ \begin{cases} y = x^2 + 6x + 10 \\ y = 1 \end{cases} \]
\begin{tasks}(4)
    \task -3
    \task -10
    \task 0
    \task 3
\end{tasks}

\begin{solution}
\textbf{Conceptual Explanation:} \\
Since the system has exactly one solution, the line $y = 1$ must touch the vertex of the parabola. We can find the $x$-coordinate of the vertex using the formula $x = -b/2a$.

\vspace{20pt}

\textbf{Calculation and Logic:} \\
For the parabola $y = x^2 + 6x + 10$: \\
$a = 1, b = 6$ \\
$x = -b/2a = -6 / 2(1) = -3$

\vspace{10pt}

Verify the $y$-value at $x = -3$: \\
$y = (-3)^2 + 6(-3) + 10$ \\
$y = 9 - 18 + 10 = 1$ \\
This confirms that the point $(-3, 1)$ is the only solution.

\vspace{25pt}

Answer: \colorbox{yellow!100}{A}
\end{solution}

\vspace{40pt}

% Question 19
% Category: Practice | Difficulty: Medium | Question Type: Multiple Choice
\question
Which of the following describes the number of real solutions to the system below?
\[ \begin{cases} y = 2x^2 + 5 \\ y = x + 1 \end{cases} \]
\begin{tasks}(4)
    \task Zero
    \task One
    \task Two
    \task Infinitely many
\end{tasks}

\begin{solution}
\textbf{Conceptual Explanation:} \\
Calculate the discriminant of the combined equation. A negative discriminant means no real solutions.

\vspace{20pt}

\textbf{Calculation and Logic:} \\
$2x^2 + 5 = x + 1$ \\
$2x^2 - x + 4 = 0$

\vspace{10pt}

$D = (-1)^2 - 4(2)(4)$ \\
$D = 1 - 32 = -31$ \\
Since $D < 0$, there are zero real solutions.

\vspace{25pt}

Answer: \colorbox{yellow!100}{A}
\end{solution}

\vspace{40pt}

% Question 20
% Category: Practice | Difficulty: Hard | Question Type: Multiple Choice
\question
A circle is centered at the origin and has a radius of 5. A line is defined by $y = 3$. At how many points do the circle and the line intersect?
\begin{tasks}(4)
    \task Zero
    \task One
    \task Two
    \task Three
\end{tasks}

\begin{solution}
\textbf{Conceptual Explanation:} \\
A circle centered at the origin with radius $r$ has a maximum height of $r$ and a minimum height of $-r$. If a horizontal line is within this range, it will intersect the circle.

\vspace{20pt}

\textbf{Calculation and Logic:} \\
The circle reaches from $y = -5$ to $y = 5$ on the vertical axis. \\
The line $y = 3$ is between the top and bottom of the circle ($|3| < 5$). \\
Therefore, the line will pass through the circle and intersect it at two distinct points.

\vspace{25pt}

Answer: \colorbox{yellow!100}{C}
\end{solution}

\vspace{40pt}

% Question 21
% Category: Practice | Difficulty: Medium | Question Type: Multiple Choice
\question
If $(x, y)$ is a solution to the system below, what is a possible value for $y$?
\[ \begin{cases} y = x^2 - 4 \\ y = -3 \end{cases} \]
\begin{tasks}(4)
    \task -4
    \task -3
    \task 1
    \task 4
\end{tasks}

\begin{solution}
\textbf{Conceptual Explanation:} \\
A system involving a constant line $y = k$ means that for any solution $(x, y)$, the $y$-value must be $k$.

\vspace{20pt}

\textbf{Calculation and Logic:} \\
Since the second equation is $y = -3$, any intersection point must have a $y$-coordinate of $-3$.

\vspace{25pt}

Answer: \colorbox{yellow!100}{B}
\end{solution}

\vspace{40pt}

% Question 22
% Category: Practice | Difficulty: Hard | Question Type: Multiple Choice
\question
The parabola $y = x^2 + 4x + k$ and the line $y = 2x + 1$ intersect at exactly one point. What is the value of $k$?
\begin{tasks}(4)
    \task 0
    \task 1
    \task 2
    \task 3
\end{tasks}

\begin{solution}
\textbf{Conceptual Explanation:} \\
Set the equations equal and use the condition that the discriminant must be zero for exactly one solution.

\vspace{20pt}

\textbf{Calculation and Logic:} \\
$x^2 + 4x + k = 2x + 1$ \\
$x^2 + 2x + (k - 1) = 0$

\vspace{10pt}

$D = (2)^2 - 4(1)(k - 1) = 0$ \\
$4 - 4k + 4 = 0$ \\
$8 - 4k = 0$ \\
$4k = 8 \Rightarrow k = 2$

\vspace{25pt}

Answer: \colorbox{yellow!100}{C}
\end{solution}

\vspace{40pt}

% Question 23
% Category: Practice | Difficulty: Medium | Question Type: Multiple Choice
\question
Find the $x$-coordinates of the intersection points of $y = x^2 - 16$ and the $x$-axis.
\begin{tasks}(4)
    \task $x = 0$
    \task $x = 4$ and $x = -4$
    \task $x = 16$ and $x = -16$
    \task $x = 8$ and $x = -8$
\end{tasks}

\begin{solution}
\textbf{Conceptual Explanation:} \\
The $x$-axis is defined by the equation $y = 0$. Finding the intersection points with the $x$-axis means solving the quadratic equation when $y$ is set to zero.

\vspace{20pt}

\textbf{Calculation and Logic:} \\
$x^2 - 16 = 0$ \\
$x^2 = 16$ \\
$x = \pm 4$

\vspace{25pt}

Answer: \colorbox{yellow!100}{B}
\end{solution}

\vspace{40pt}

% Question 24
% Category: Practice | Difficulty: Hard | Question Type: Multiple Choice
\question
If $y = x^2$ and $y = mx - 9$ intersect at only one point and $m > 0$, what is the point of intersection?
\begin{tasks}(4)
    \task $(3, 9)$
    \task $(6, 36)$
    \task $(0, 0)$
    \task $(9, 81)$
\end{tasks}

\begin{solution}
\textbf{Conceptual Explanation:} \\
Find $m$ using the discriminant first, then solve for the intersection point.

\vspace{20pt}

\textbf{Calculation and Logic:} \\
$x^2 - mx + 9 = 0$ \\
$D = m^2 - 4(1)(9) = 0 \Rightarrow m = 6$ \\
Now substitute $m = 6$: \\
$x^2 - 6x + 9 = 0$ \\
$(x - 3)^2 = 0 \Rightarrow x = 3$

\vspace{10pt}

Find $y$ at $x = 3$: \\
$y = (3)^2 = 9$ \\
The point is $(3, 9)$.

\vspace{25pt}

Answer: \colorbox{yellow!100}{A}
\end{solution}

\vspace{40pt}

% Question 25
% Category: Practice | Difficulty: Medium | Question Type: Multiple Choice
\question

Based on the graph above, which of the following is the equation of the line?
\begin{tasks}(2)
    \task $y = 3x + 2$
    \task $y = x + 2$
    \task $y = 3x$
    \task $y = 2x + 3$
\end{tasks}

\begin{solution}
\textbf{Conceptual Explanation:} \\
Identify two points on the line from the graph and use the slope-intercept form $y = mx + b$.

\vspace{20pt}

\textbf{Calculation and Logic:} \\
Points on the line: $(0, 2)$ and $(3, 11)$. \\
$y$-intercept ($b$) = 2. \\
Slope ($m$) = $(11 - 2) / (3 - 0) = 9 / 3 = 3$. \\
Equation: $y = 3x + 2$.

\vspace{25pt}

Answer: \colorbox{yellow!100}{A}
\end{solution}

% --- PART 2: Advanced Linear-Quadratic Systems (26-50) ---

% Question 26
% Category: Practice | Difficulty: Hard | Question Type: Multiple Choice
\question
The parabola $y = x^2 - 10x + 21$ and the line $y = -x + k$ intersect at exactly two points. If one of the intersection points is $(4, -3)$, what is the $x$-coordinate of the other intersection point?
\begin{tasks}(4)
    \task 5
    \task 3
    \task 7
    \task 2
\end{tasks}

\begin{solution}
\textbf{Conceptual Explanation:} \\
Since $(4, -3)$ is an intersection point, it must satisfy both equations. We first find the value of $k$ using the linear equation, then set the quadratic and linear equations equal to find the second solution.

\vspace{20pt}

\textbf{Calculation and Logic:} \\
Substitute $(4, -3)$ into $y = -x + k$: \\
$-3 = -4 + k \Rightarrow k = 1$.

\vspace{10pt}

Set the equations equal: \\
$x^2 - 10x + 21 = -x + 1$ \\
$x^2 - 9x + 20 = 0$

\vspace{10pt}

Factor the quadratic: \\
$(x - 4)(x - 5) = 0$ \\
The $x$-coordinates are 4 and 5.

\vspace{10pt}

The other $x$-coordinate is 5.

\vspace{25pt}

Answer: \colorbox{yellow!100}{A}
\end{solution}

\vspace{40pt}

% Question 27
% Category: Practice | Difficulty: Medium | Question Type: Multiple Choice
\question
Which of the following describes the number of real solutions to the system below?
\[ \begin{cases} y = (x - 4)^2 + 5 \\ y = x \end{cases} \]
\begin{tasks}(4)
    \task Zero
    \task One
    \task Two
    \task Infinitely many
\end{tasks}

\begin{solution}
\textbf{Conceptual Explanation:} \\
Expanding the vertex form and setting it equal to the linear equation allows us to use the discriminant to determine the number of intersection points.

\vspace{20pt}

\textbf{Calculation and Logic:} \\
$(x - 4)^2 + 5 = x$ \\
$x^2 - 8x + 16 + 5 = x$ \\
$x^2 - 9x + 21 = 0$

\vspace{10pt}

Calculate the discriminant ($b^2 - 4ac$): \\
$D = (-9)^2 - 4(1)(21)$ \\
$D = 81 - 84 = -3$

\vspace{10pt}

Since the discriminant is negative, there are zero real solutions.

\vspace{25pt}

Answer: \colorbox{yellow!100}{A}
\end{solution}

\vspace{40pt}

% Question 28
% Category: Practice | Difficulty: Hard | Question Type: Multiple Choice
\question
The line $y = 3$ intersects the circle $(x - 1)^2 + (y + 1)^2 = 25$ at two points. What is the $x$-coordinate of the point with the larger $x$-value?
\begin{tasks}(4)
    \task 1
    \task 4
    \task 5
    \task 6
\end{tasks}

\begin{solution}
\textbf{Conceptual Explanation:} \\
Substitute the $y$-value into the circle's equation and solve for $x$. This identifies where the horizontal line crosses the boundary of the circle.

\vspace{20pt}

\textbf{Calculation and Logic:} \\
Substitute $y = 3$: \\
$(x - 1)^2 + (3 + 1)^2 = 25$ \\
$(x - 1)^2 + 4^2 = 25$ \\
$(x - 1)^2 + 16 = 25$

\vspace{10pt}

Isolate $(x - 1)^2$: \\
$(x - 1)^2 = 9$ \\
$x - 1 = \pm 3$

\vspace{10pt}

Possible $x$-values: \\
$x = 1 + 3 = 4$ \\
$x = 1 - 3 = -2$

\vspace{10pt}

The larger $x$-value is 4.

\vspace{25pt}

Answer: \colorbox{yellow!100}{B}
\end{solution}

\vspace{40pt}

% Question 29
% Category: Practice | Difficulty: Medium | Question Type: Multiple Choice
\question
For what value of $c$ will the parabola $y = x^2 + c$ pass through the point $(2, 10)$?
\begin{tasks}(4)
    \task 2
    \task 4
    \task 6
    \task 8
\end{tasks}

\begin{solution}
\textbf{Conceptual Explanation:} \\
If a point lies on the graph of a function, substituting the point's coordinates into the function's equation must result in a true statement.

\vspace{20pt}

\textbf{Calculation and Logic:} \\
Substitute $x = 2$ and $y = 10$ into $y = x^2 + c$: \\
$10 = 2^2 + c$ \\
$10 = 4 + c$ \\
$c = 6$

\vspace{25pt}

Answer: \colorbox{yellow!100}{C}
\end{solution}

\vspace{40pt}

% Question 30
% Category: Practice | Difficulty: Hard | Question Type: Multiple Choice
\question
A system of equations is given by $y = x^2 + kx + 4$ and $y = 0$. If the system has exactly one solution, what is one possible value for $k$?
\begin{tasks}(4)
    \task 2
    \task 4
    \task 8
    \task 16
\end{tasks}

\begin{solution}
\textbf{Conceptual Explanation:} \\
Exactly one solution means the discriminant of $x^2 + kx + 4 = 0$ must be zero. This signifies that the parabola's vertex is on the $x$-axis.

\vspace{20pt}

\textbf{Calculation and Logic:} \\
$D = k^2 - 4(1)(4) = 0$ \\
$k^2 - 16 = 0$ \\
$k^2 = 16$ \\
$k = \pm 4$

\vspace{10pt}

Among the choices, 4 is the correct answer.

\vspace{25pt}

Answer: \colorbox{yellow!100}{B}
\end{solution}

\vspace{40pt}

% Question 31
% Category: Practice | Difficulty: Hard | Question Type: Multiple Choice
\question
If $y = 2x^2 - 4x + 7$ and $y = 2x + 3$, what is the distance between the two points of intersection?
\begin{tasks}(4)
    \task $\sqrt{5}$
    \task $2\sqrt{5}$
    \task $2\sqrt{2}$
    \task 4
\end{tasks}

\begin{solution}
\textbf{Conceptual Explanation:} \\
First, find the $x$-coordinates of the intersection points. Then find the corresponding $y$-coordinates and use the distance formula: $d = \sqrt{(x_2-x_1)^2 + (y_2-y_1)^2}$.

\vspace{20pt}

\textbf{Calculation and Logic:} \\
$2x^2 - 4x + 7 = 2x + 3$ \\
$2x^2 - 6x + 4 = 0 \Rightarrow x^2 - 3x + 2 = 0$ \\
$(x - 2)(x - 1) = 0 \Rightarrow x = 2, x = 1$

\vspace{10pt}

Find points: \\
If $x = 1, y = 2(1) + 3 = 5 \Rightarrow (1, 5)$ \\
If $x = 2, y = 2(2) + 3 = 7 \Rightarrow (2, 7)$

\vspace{10pt}

Calculate distance: \\
$d = \sqrt{(2-1)^2 + (7-5)^2} = \sqrt{1^2 + 2^2} = \sqrt{5}$

\vspace{25pt}

Answer: \colorbox{yellow!100}{A}
\end{solution}

\vspace{40pt}

% Question 32
% Category: Practice | Difficulty: Medium | Question Type: Multiple Choice
\question
How many points of intersection are there between $x^2 + y^2 = 9$ and $y = x^2 + 10$?
\begin{tasks}(4)
    \task Zero
    \task One
    \task Two
    \task Four
\end{tasks}

\begin{solution}
\textbf{Conceptual Explanation:} \\
Compare the vertical ranges of the two graphs. The circle is centered at the origin with a radius of 3. The parabola opens upward from a vertex at $(0, 10)$.

\vspace{20pt}

\textbf{Calculation and Logic:} \\
The highest point on the circle is $y = 3$. \\
The lowest point on the parabola is $y = 10$. \\
Since the entire parabola is located 7 units above the circle, they cannot intersect.

\vspace{25pt}

Answer: \colorbox{yellow!100}{A}
\end{solution}

\vspace{40pt}

% Question 33
% Category: Practice | Difficulty: Hard | Question Type: Multiple Choice
\question
If $y = x^2 - 6x + 10$ and $y = mx + 1$, what value of $m$ ensures the line passes through the vertex of the parabola?
\begin{tasks}(4)
    \task 0
    \task 1
    \task 2
    \task -2
\end{tasks}

\begin{solution}
\textbf{Conceptual Explanation:} \\
First find the vertex of the parabola using $x = -b/2a$. Then substitute the vertex coordinates into the linear equation to solve for $m$.

\vspace{20pt}

\textbf{Calculation and Logic:} \\
Vertex $x = -(-6)/2 = 3$. \\
Vertex $y = 3^2 - 6(3) + 10 = 1$. \\
Vertex is $(3, 1)$.

\vspace{10pt}

Substitute into $y = mx + 1$: \\
$1 = m(3) + 1$ \\
$0 = 3m \Rightarrow m = 0$

\vspace{25pt}

Answer: \colorbox{yellow!100}{A}
\end{solution}

\vspace{40pt}

% Question 34
% Category: Practice | Difficulty: Medium | Question Type: Multiple Choice
\question
Which of the following lines is tangent to the parabola $y = x^2$ at the point $(0, 0)$?
\begin{tasks}(4)
    \task $y = x$
    \task $y = -x$
    \task $y = 0$
    \task $x = 0$
\end{tasks}

\begin{solution}
\textbf{Conceptual Explanation:} \\
The vertex of $y = x^2$ is $(0, 0)$. The tangent line at the vertex of a vertical parabola is always a horizontal line.

\vspace{20pt}

\textbf{Calculation and Logic:} \\
The horizontal line passing through $(0, 0)$ has the equation $y = 0$.

\vspace{25pt}

Answer: \colorbox{yellow!100}{C}
\end{solution}

\vspace{40pt}

% Question 35
% Category: Practice | Difficulty: Hard | Question Type: Multiple Choice
\question

The system consists of $x^2 + y^2 = 20$ and $y = -2x$. What is the $y$-coordinate of the solution in the second quadrant?
\begin{tasks}(4)
    \task 2
    \task 4
    \task -4
    \task -2
\end{tasks}

\begin{solution}
\textbf{Conceptual Explanation:} \\
Substitute the linear relationship for $y$ into the circle equation to find $x$, then calculate $y$. Recall that in the second quadrant, $x < 0$ and $y > 0$.

\vspace{20pt}

\textbf{Calculation and Logic:} \\
$x^2 + (-2x)^2 = 20$ \\
$x^2 + 4x^2 = 20 \Rightarrow 5x^2 = 20$ \\
$x^2 = 4 \Rightarrow x = \pm 2$.

\vspace{10pt}

For the second quadrant, $x = -2$. \\
$y = -2(-2) = 4$.

\vspace{25pt}

Answer: \colorbox{yellow!100}{B}
\end{solution}

\vspace{40pt}

% Question 36
% Category: Practice | Difficulty: Medium | Question Type: Multiple Choice
\question
At how many points does $y = x^2 - 4x + 4$ intersect the $x$-axis?
\begin{tasks}(4)
    \task 0
    \task 1
    \task 2
    \task 4
\end{tasks}

\begin{solution}
\textbf{Conceptual Explanation:} \\
The $x$-axis is the line $y = 0$. Factor the quadratic to see if it has a repeated root.

\vspace{20pt}

\textbf{Calculation and Logic:} \\
$x^2 - 4x + 4 = 0$ \\
$(x - 2)^2 = 0$ \\
$x = 2$ (one distinct root).

\vspace{10pt}

There is exactly one point of intersection.

\vspace{25pt}

Answer: \colorbox{yellow!100}{B}
\end{solution}

\vspace{40pt}

% Question 37
% Category: Practice | Difficulty: Hard | Question Type: Multiple Choice
\question
The parabola $y = ax^2 + bx + c$ has vertex $(2, 3)$. If the line $y = 3$ is tangent to the parabola, which of the following could be the value of $a$?
\begin{tasks}(4)
    \task Any positive number
    \task Any negative number
    \task Any non-zero number
    \task Only $a = 1$
\end{tasks}

\begin{solution}
\textbf{Conceptual Explanation:} \\
A horizontal line $y = k$ is tangent to a parabola at its vertex. This remains true regardless of how "steep" (the value of $a$) the parabola is, as long as it opens up or down.

\vspace{20pt}

\textbf{Calculation and Logic:} \\
The vertex is $(2, 3)$. The line $y = 3$ is horizontal and passes through the vertex. Thus, any parabola with this vertex will be tangent to the line.

\vspace{25pt}

Answer: \colorbox{yellow!100}{C}
\end{solution}

\vspace{40pt}

% Question 38
% Category: Practice | Difficulty: Medium | Question Type: Multiple Choice
\question
If $y = x^2 + 6x + 5$ and $y = 5$, what is the sum of the $x$-coordinates of the intersection points?
\begin{tasks}(4)
    \task 6
    \task -6
    \task 0
    \task 5
\end{tasks}

\begin{solution}
\textbf{Conceptual Explanation:} \\
Set the quadratic equal to the constant and use the sum of roots formula $-b/a$.

\vspace{20pt}

\textbf{Calculation and Logic:} \\
$x^2 + 6x + 5 = 5$ \\
$x^2 + 6x = 0$ \\
$a = 1, b = 6$. \\
Sum $= -6/1 = -6$.

\vspace{25pt}

Answer: \colorbox{yellow!100}{B}
\end{solution}

\vspace{40pt}

% Question 39
% Category: Practice | Difficulty: Hard | Question Type: Multiple Choice
\question
What is the $y$-intercept of the line that is tangent to $y = x^2$ at the point $(2, 4)$? (Assume the slope is 4).
\begin{tasks}(4)
    \task 0
    \task 4
    \task -4
    \task 8
\end{tasks}

\begin{solution}
\textbf{Conceptual Explanation:} \\
Use the point-slope form $y - y_1 = m(x - x_1)$ with $m = 4$ and point $(2, 4)$, then solve for $y$ when $x = 0$.

\vspace{20pt}

\textbf{Calculation and Logic:} \\
$y - 4 = 4(x - 2)$ \\
$y - 4 = 4x - 8$ \\
$y = 4x - 4$ \\
The $y$-intercept is $-4$.

\vspace{25pt}

Answer: \colorbox{yellow!100}{C}
\end{solution}

\vspace{40pt}

% Question 40
% Category: Practice | Difficulty: Medium | Question Type: Multiple Choice
\question
How many real solutions are there for the system $y = x^2 + 1$ and $y = x^2 + 2$?
\begin{tasks}(4)
    \task Zero
    \task One
    \task Two
    \task Infinitely many
\end{tasks}

\begin{solution}
\textbf{Conceptual Explanation:} \\
These are two parabolas with the same shape but different vertical shifts.

\vspace{20pt}

\textbf{Calculation and Logic:} \\
$x^2 + 1 = x^2 + 2$ \\
$1 = 2$ (impossible). \\
There are no real solutions.

\vspace{25pt}

Answer: \colorbox{yellow!100}{A}
\end{solution}

\vspace{40pt}

% Question 41
% Category: Practice | Difficulty: Hard | Question Type: Multiple Choice
\question
If the line $y = 2x + k$ is tangent to the parabola $y = x^2$, what is the value of $k$?
\begin{tasks}(4)
    \task 1
    \task -1
    \task 0
    \task 2
\end{tasks}

\begin{solution}
\textbf{Conceptual Explanation:} \\
Set the equations equal and set the discriminant of the resulting quadratic to zero.

\vspace{20pt}

\textbf{Calculation and Logic:} \\
$x^2 = 2x + k$ \\
$x^2 - 2x - k = 0$ \\
$D = (-2)^2 - 4(1)(-k) = 0$ \\
$4 + 4k = 0 \Rightarrow k = -1$.

\vspace{25pt}

Answer: \colorbox{yellow!100}{B}
\end{solution}

\vspace{40pt}

% Question 42
% Category: Practice | Difficulty: Medium | Question Type: Multiple Choice
\question
Which of the following points is a solution to $y = x^2 - x - 6$ and $y = 0$?
\begin{tasks}(4)
    \task $(3, 0)$
    \task $(-3, 0)$
    \task $(2, 0)$
    \task $(6, 0)$
\end{tasks}

\begin{solution}
\textbf{Conceptual Explanation:} \\
Find the $x$-intercepts of the parabola by factoring.

\vspace{20pt}

\textbf{Calculation and Logic:} \\
$x^2 - x - 6 = 0$ \\
$(x - 3)(x + 2) = 0$ \\
$x = 3, x = -2$. \\
$(3, 0)$ is the solution provided in the choices.

\vspace{25pt}

Answer: \colorbox{yellow!100}{A}
\end{solution}

\vspace{40pt}

% Question 43
% Category: Practice | Difficulty: Hard | Question Type: Multiple Choice
\question
The line $y = x + 1$ and the parabola $y = x^2 - x - 2$ intersect at two points. What is the $y$-coordinate of the point with the larger $y$-value?
\begin{tasks}(4)
    \task 4
    \task 3
    \task 0
    \task -2
\end{tasks}

\begin{solution}
\textbf{Conceptual Explanation:} \\
Solve for $x$ by setting the equations equal, find the corresponding $y$-values, and pick the larger one.

\vspace{20pt}

\textbf{Calculation and Logic:} \\
$x^2 - x - 2 = x + 1$ \\
$x^2 - 2x - 3 = 0 \Rightarrow (x - 3)(x + 1) = 0$ \\
$x = 3, x = -1$.

\vspace{10pt}

If $x = 3, y = 3 + 1 = 4$. \\
If $x = -1, y = -1 + 1 = 0$. \\
The larger $y$-value is 4.

\vspace{25pt}

Answer: \colorbox{yellow!100}{A}
\end{solution}

\vspace{40pt}

% Question 44
% Category: Practice | Difficulty: Medium | Question Type: Multiple Choice
\question
The system $x^2 + y^2 = 1$ and $y = 2$ has how many real solutions?
\begin{tasks}(4)
    \task 0
    \task 1
    \task 2
    \task Infinitely many
\end{tasks}

\begin{solution}
\textbf{Conceptual Explanation:} \\
The circle has a radius of 1. The line $y = 2$ is outside the boundary of the circle.

\vspace{20pt}

\textbf{Calculation and Logic:} \\
$x^2 + 2^2 = 1 \Rightarrow x^2 + 4 = 1 \Rightarrow x^2 = -3$. \\
No real solutions.

\vspace{25pt}

Answer: \colorbox{yellow!100}{A}
\end{solution}

\vspace{40pt}

% Question 45
% Category: Practice | Difficulty: Hard | Question Type: Multiple Choice
\question

How many total real solutions could a system consisting of a circle and a parabola have?
\begin{tasks}(4)
    \task At most 2
    \task At most 4
    \task Exactly 2
    \task Infinitely many
\end{tasks}

\begin{solution}
\textbf{Conceptual Explanation:} \\
A circle (quadratic in $x$ and $y$) and a parabola (quadratic in $x$) can cross at several points depending on their orientation.

\vspace{20pt}

\textbf{Calculation and Logic:} \\
By sketching, we can find cases where the parabola dips into and out of the circle, creating up to 4 points of intersection.

\vspace{25pt}

Answer: \colorbox{yellow!100}{B}
\end{solution}

\vspace{40pt}

% Question 46
% Category: Practice | Difficulty: Medium | Question Type: Multiple Choice
\question
The parabola $y = x^2 - 2x + 1$ is tangent to the $x$-axis. What is the $x$-coordinate of the intersection?
\begin{tasks}(4)
    \task 0
    \task 1
    \task 2
    \task -1
\end{tasks}

\begin{solution}
\textbf{Conceptual Explanation:} \\
The intersection with the $x$-axis occurs at the vertex since the parabola is tangent to it.

\vspace{20pt}

\textbf{Calculation and Logic:} \\
$x^2 - 2x + 1 = 0 \Rightarrow (x - 1)^2 = 0 \Rightarrow x = 1$.

\vspace{25pt}

Answer: \colorbox{yellow!100}{B}
\end{solution}

\vspace{40pt}

% Question 47
% Category: Practice | Difficulty: Hard | Question Type: Multiple Choice
\question
If $y = x^2$ and $y = 2x + 3$, what is the sum of the $y$-coordinates of the solutions?
\begin{tasks}(4)
    \task 10
    \task 2
    \task 9
    \task 4
\end{tasks}

\begin{solution}
\textbf{Conceptual Explanation:} \\
Find both $x$ values, find their corresponding $y$ values, then add the $y$ values together.

\vspace{20pt}

\textbf{Calculation and Logic:} \\
$x^2 = 2x + 3 \Rightarrow x^2 - 2x - 3 = 0 \Rightarrow (x - 3)(x + 1) = 0$ \\
$x = 3, x = -1$.

\vspace{10pt}

If $x = 3, y = 3^2 = 9$. \\
If $x = -1, y = (-1)^2 = 1$. \\
Sum of $y = 9 + 1 = 10$.

\vspace{25pt}

Answer: \colorbox{yellow!100}{A}
\end{solution}

\vspace{40pt}

% Question 48
% Category: Practice | Difficulty: Medium | Question Type: Multiple Choice
\question
Does the system $y = x^2 + 5x + 6$ and $y = x + 1$ have any real solutions?
\begin{tasks}(4)
    \task Yes, two solutions
    \task Yes, one solution
    \task No real solutions
    \task Infinitely many
\end{tasks}

\begin{solution}
\textbf{Conceptual Explanation:} \\
Check the discriminant of the combined equation.

\vspace{20pt}

\textbf{Calculation and Logic:} \\
$x^2 + 5x + 6 = x + 1 \Rightarrow x^2 + 4x + 5 = 0$ \\
$D = 4^2 - 4(1)(5) = 16 - 20 = -4$. \\
No real solutions.

\vspace{25pt}

Answer: \colorbox{yellow!100}{C}
\end{solution}

\vspace{40pt}

% Question 49
% Category: Practice | Difficulty: Hard | Question Type: Multiple Choice
\question
The line $y = 4$ and the parabola $y = (x - h)^2 + k$ intersect at exactly one point. What is the value of $k$?
\begin{tasks}(4)
    \task $h$
    \task 0
    \task 4
    \task Cannot be determined
\end{tasks}

\begin{solution}
\textbf{Conceptual Explanation:} \\
For a horizontal line to be tangent to a parabola in vertex form, the $y$-coordinate of the vertex ($k$) must equal the $y$-coordinate of the line.

\vspace{20pt}

\textbf{Calculation and Logic:} \\
Vertex $y = k$. Line $y = 4$. \\
$k = 4$.

\vspace{25pt}

Answer: \colorbox{yellow!100}{C}
\end{solution}

\vspace{40pt}

% Question 50
% Category: Practice | Difficulty: Hard | Question Type: Multiple Choice
\question
The system $y = x^2 + x$ and $y = c$ has no real solutions. What is the maximum possible integer value for $c$?
\begin{tasks}(4)
    \task 0
    \task -1
    \task 1
    \task -2
\end{tasks}

\begin{solution}
\textbf{Conceptual Explanation:} \\
Find the minimum $y$-value of the parabola (the vertex). Any value of $c$ less than this minimum will result in no solutions.

\vspace{20pt}

\textbf{Calculation and Logic:} \\
$x$-vertex $= -1/2 = -0.5$. \\
$y$-vertex $= (-0.5)^2 + (-0.5) = 0.25 - 0.5 = -0.25$. \\
The minimum $y$-value is $-0.25$. \\
For no solutions, $c < -0.25$. \\
The largest integer less than $-0.25$ is $-1$.

\vspace{25pt}

Answer: \colorbox{yellow!100}{B}
\end{solution}

% =============================================================
% Question End
% =============================================================

\end{questions}
\begin{center}
\rule{.5\textwidth}{1pt}
\end{center}
\end{document}
